\section{Метод}\label{sec3}

На рис.~\ref{fig:scheme} представлена схема декодирования изображений. Процесс строится из двух основных этапов: контрастивное обучение комбинированных представлений фМРТ--ЭЭГ сигналов и реконструкция стимулов.

\begin{figure}[h]
    \centering
    \includegraphics[width=1.\textwidth]{figure_1.pdf}
    \caption{Схема декодирования изображений на основе сигналов фМРТ--ЭЭГ. Процесс декодирования реализуется в два этапа. На первом этапе проводится контрастное обучение энкодера сигналов фМРТ--ЭЭГ, в ходе которого представления мозговых сигналов выравниваются с замороженными эмбеддингами изображений, полученными с помощью модели CLIP. На втором этапе осуществляется двухстадийная генерация реконструкций: сначала обучается диффузионная априорная модель для уточнения эмбеддингов сигналов мозга, а затем предобученная диффузионная модель, обусловленная на CLIP-эмбеддинги, генерирует реконструкции изображений}
    \label{fig:scheme}
\end{figure}

\subsection{Контрастивное обучение комбинированных представлений фМРТ--ЭЭГ сигналов}\label{sec3:contrastive}

\paragraph{Энкодер сигналов фМРТ--ЭЭГ.}
Архитектура энкодера состоит из трех модулей: энкодера фМРТ, энкодера ЭЭГ и модуля объединения. Для обеспечения персонализированного представления мозговой активности каждый модуль включает индивидуализирующие компоненты, позволяющие явно моделировать межсубъектные различия~--- критически важный аспект при работе с нейросигналами. Архитектуры энкодеров фМРТ и ЭЭГ изображены на рис.~\ref{fig:encoders}.

\paragraph{Предобработка ЭЭГ.}
На этапе предобработки решается задача восстановления пропусков в записях ЭЭГ: сигналы незарегистрированных электродов восстанавливаются методами пространственной интерполяции на основе данных ближайших доступных каналов и известной геометрии расположения датчиков.

Полный набор электродов системы регистрации включает $K$ каналов. Каждому каналу $\{c_k\}_{k=1}^K$ соответствуют фиксированные пространственные координаты $\mathbf{p}_k \in \bbR^3$ на поверхности скальпа. Множество фактически зарегистрированных каналов для $i$-й записи обозначается как $\mathcal{C}_i \subseteq \{c_1, \ldots, c_K\},~ i=1,\ldots, N$. Состав отсутствующих каналов варьируется между записями, вследствие чего использование исключительно перекрывающегося подмножества каналов приводит к существенной потере информации. Архитектура энкодера ЭЭГ $\mathcal{E}_e$ предполагает фиксированную размерность входного тензора, что необходимо для обработки элементов обучающей выборки батчами. Данное требование обусловливает необходимость восстановления сигналов отсутствующих каналов до полного набора.

Для упрощения нотации далее рассматривается единичная запись ЭЭГ $\mathbf{E}$, при этом индекс $i$ объекта выборки опускается. Обозначим через $\mathcal{C} \subset \{c_1, \dots, c_K\}$~--- подмножество каналов, для которых доступны временные ряды $\mathbf{e}_k \in \bbR^T$. Таким образом, наблюдаемая часть записи задается как $\mathbf{E} = \{\mathbf{e}_k\}_{k \in \mathcal{C}}$.
Для унификации формата данных каждому отсутствующему каналу $c_k \notin \mathcal{C}$ назначается восстановленный сигнал $\widehat{\mathbf{e}}_k \in \bbR^T$, рассчитываемый как взвешенное среднее значений сигналов $n$ пространственно ближайших зарегистрированных электродов:
\begin{equation}\label{eq:recover}
\widehat{\mathbf{e}}_k = \sum_{j \in \mathcal{S}_k} w_{kj}\, \mathbf{e}_j, \qquad \sum_{j \in \mathcal{S}_k} w_{kj} = 1,
\end{equation}
где $\mathbf{e}_j \in \bbR^T$~--- временной ряд канала $c_j$; $\mathcal{S}_k$~--- множество индексов $n$ ближайших к $c_k$ доступных электродов, определяемое в соответствии с выбранной пространственной метрикой; $w_{kj} \geq 0$~--- неотрицательные нормированные веса.
В исследовании рассматриваются четыре стратегии определения весов $\{w_{kj}\}$ и формирования множества соседей $\mathcal{S}_k$:
\begin{itemize}
    \item \textbf{Zero} (нулевое заполнение): базовый метод, при котором отсутствующий сигнал заменяется нулевым вектором:
    \begin{equation*}
    \widehat{\mathbf{e}}_k = \mathbf{0} \in \bbR^T.
    \end{equation*}
    \item \textbf{NN} (метод одного ближайшего соседа): сигнал отсутствующего канала заменяется сигналом ближайшего зарегистрированного электрода:
    \begin{equation*}
    \widehat{\mathbf{e}}_k = \mathbf{e}_{j^*}, \qquad j^* = \argmin_{j: c_j \in \mathcal{C}} \rho_{\text{eucl}}(c_k, c_j), \quad \rho_{\text{eucl}}(c_k, c_j) = \|\mathbf{p}_k - \mathbf{p}_j\|_2.
    \end{equation*}
    что соответствует $n = 1$ и $w_{kj^*} = 1$ в~\eqref{eq:recover}.
    \item \textbf{KNN-Euclidean}: взвешенное усреднение по $n$ ближайшим соседям в евклидовой метрике; веса определяются как величины, обратно пропорциональные расстоянию между электродами:
    \begin{equation*}
    w_{kj} = \frac{1 / \rho_{\text{eucl}}(c_k, c_j)}{\displaystyle\sum_{l \in \mathcal{S}_k} 1 / \rho_{\text{eucl}}(c_k, c_l)}, \qquad \rho_{\text{eucl}}(c_k, c_j) = \|\mathbf{p}_k - \mathbf{p}_j\|_2.
    \end{equation*}
    \item \textbf{KNN-Spherical}: взвешенное усреднение по $n$ ближайшим соседям в сферической метрике; процедура аналогична предыдущей, однако евклидово расстояние заменяется геодезическим расстоянием на поверхности скальпа, аппроксимируемой сферой радиуса $R$:
    \begin{equation*}
    \rho_{\text{sph}}(c_k, c_j) = R \cdot \arccos\!\left( \frac{\langle \mathbf{p}_k, \mathbf{p}_j \rangle}{\|\mathbf{p}_k\|_2 \, \|\mathbf{p}_j\|_2} \right),
    \end{equation*}
    при этом веса задаются как $w_{kj} \propto 1 / \rho_{\text{sph}}(c_k, c_j)$ с последующей нормировкой $\sum_{j \in \mathcal{S}_k} w_{kj} = 1$.
\end{itemize}
Применение процедуры~\eqref{eq:recover} ко всем отсутствующим каналам приводит исходные данные ЭЭГ к единому матричному представлению $\mathbf{E} \in \bbR^{K \times T}$ с фиксированным числом каналов $K$. Данный шаг обеспечивает совместимость размерности входного тензора с архитектурой энкодера $\mathcal{E}_e$ для последующей обработки батчами.

\paragraph{Энкодер ЭЭГ.}
Энкодер ЭЭГ интегрирует индивидуальные обучаемые эмбеддинги пациентов с исходными многоканальными данными посредством конкатенации, формируя единое векторное представление. Архитектура энкодера включает последовательность трех сверточных слоев: одномерную свертку по временному измерению, одномерную свертку по каналам (пространству) и свертку для сжатия в вектор. Полученный эмбеддинг пропускается через MLP-голову, формируя итоговое представление. Формально, для батча данных ЭЭГ $\mathbf{E} \in \mathbb{R}^{B \times K \times T}$ энкодер $\mathcal{E}_{e}$ генерирует эмбеддинги:
\begin{equation*}
\mathbf{Z}_{e} = \mathcal{E}_{e}(\mathbf{E}) \in \mathbb{R}^{B \times d_{e}},
\end{equation*}
где $d_{e}$~--- размерность латентного пространства ЭЭГ.

\paragraph{Предобработка фМРТ.}
Из всего объема $X \times Y \times Z$ вокселей лишь относительно небольшая часть несет информацию о реакции мозга на визуальный стимул: значительную долю занимают фоновые воксели, значения которых либо близки к нулю, либо стабильны во времени и не зависят от стимула. Прямая подача всего тензора в обучаемую модель приводит к расточительной параметризации и снижает отношение сигнал/шум. Поэтому на этапе предобработки строится индивидуальная для каждого испытуемого бинарная маска, выделяющая только информативные воксели.

Пусть $\mathbf{F}(t, i, j, k)$~--- значение вокселя с координатами $(i, j, k)$ в момент $t$ временного ряда фМРТ длиной $T_{f}$. Для каждого вокселя вычисляется суммарная абсолютная разность значений между соседними кадрами:
\begin{equation*}\label{eq:mask_score}
m_{ijk} = \sum_{t=2}^{T_{f}} \left| \mathbf{F}(t, i, j, k) - \mathbf{F}(t-1, i, j, k) \right|.
\end{equation*}
Величина $m_{ijk}$ интерпретируется как накопленная во времени вариация сигнала в данном вокселе: для фоновых вокселей $m_{ijk}$ близка к нулю, для активно реагирующих на стимул~--- существенно положительна. К одномерному распределению значений $\{m_{ijk}\}$ применяется метод Оцу~\cite{otsu1975threshold} для автоматического выбора порога, разделяющего фоновые и информативные воксели. Метод перебирает все возможные значения порога $\theta$; при каждом $\theta$ воксели делятся на два класса $C_0(\theta) = \{(i,j,k) : m_{ijk} \le \theta\}$ и $C_1(\theta) = \{(i,j,k) : m_{ijk} > \theta\}$ с долями $\omega_0(\theta), \omega_1(\theta) = 1 - \omega_0(\theta)$ и средними $\mu_0(\theta), \mu_1(\theta)$. Межклассовая дисперсия имеет вид
\begin{equation*}\label{eq:otsu_var}
\sigma^2(\theta) = \omega_0(\theta)\, \omega_1(\theta)\, \big[\mu_0(\theta) - \mu_1(\theta)\big]^2,
\end{equation*}
а оптимальный порог определяется как точка ее максимума:
\begin{equation*}\label{eq:otsu_threshold}
\theta^* = \argmax_{\theta} \sigma^2(\theta).
\end{equation*}
По найденному порогу строится бинарная маска
\begin{equation}\label{eq:mask}
M_{ijk} = \mathds{1}\{m_{ijk} > \theta^*\}, \qquad M \in \{0, 1\}^{X \times Y \times Z},
\end{equation}
персональная для каждого испытуемого; число информативных вокселей $d_{\text{mask}} = \sum_{ijk} M_{ijk}$ зависит от анатомии и индивидуальной динамики BOLD-сигнала. Для каждого фМРТ-скана $\mathbf{F}$ значения вокселей, для которых $M_{ijk} = 1$, собираются в плотный вектор размерности $d_{\text{mask}}$.

\paragraph{Энкодер фМРТ.}
Энкодер фМРТ учитывает индивидуальные особенности пациентов за счет персонализированной линейной проекции. Вектор информативных вокселей, отобранных по бинарной маске $M$, пропускается через индивидуальный для каждого испытуемого линейный слой, проецирующий признаки в общее латентное пространство размерности 4096. Завершающим этапом является обработка эмбеддингов через MLP с остаточными соединениями. Формально, для батча данных фМРТ $\mathbf{F} \in \mathbb{R}^{B \times X \times Y \times Z}$ энкодер $\mathcal{E}_{f}$ преобразует исходный тензор в латентное представление:
\begin{equation*}
\mathbf{Z}_{f} = \mathcal{E}_{f}(\mathbf{F}) \in \mathbb{R}^{B \times d_{f}},
\end{equation*}
где $d_{f}$~--- размерность латентного пространства фМРТ.

\paragraph{Модуль объединения.}
Полученные эмбеддинги $\mathbf{Z}_e$ и $\mathbf{Z}_f$ подаются на вход обучаемого модуля объединения $\mathcal{F}$, который интегрирует высокочастотную информацию из ЭЭГ и пространственно детализированные данные фМРТ, формируя совместные эмбеддинги:
\begin{equation*}
  \mathbf{Z}_c = \mathcal{F}(\mathbf{Z}_f, \mathbf{Z}_e) \in \mathbb{R}^{B \times d}.
\end{equation*}
В данной работе модуль объединения реализован как конкатенация эмбеддингов с последующей обработкой через MLP для извлечения совместных представлений. В качестве альтернативной реализации $\mathcal{F}$ также рассматривалось скользящее усреднение эмбеддингов модальностей вида
\begin{equation*}\label{eq:ewma}
\mathbf{Z}_c = \alpha\, \mathbf{Z}_f + (1 - \alpha)\, \mathbf{Z}_e, \qquad \alpha \in [0, 1],
\end{equation*}
где $\alpha$~--- обучаемый скалярный параметр, балансирующий вклад фМРТ- и ЭЭГ-модальностей. Сравнительный анализ показал, что MLP более эффективно выделяет совместные признаки из пространств фМРТ--ЭЭГ, поэтому именно эта конфигурация принята в качестве основной.

\begin{figure}[h]
    \centering
    \includegraphics[width=0.9\textwidth]{figure_2.pdf}
    \caption{Архитектуры кодировщиков фМРТ и ЭЭГ. \textit{Слева:} энкодер ЭЭГ интегрирует индивидуальные обучаемые эмбеддинги пациентов с исходными многоканальными данными посредством конкатенации, формируя единое векторное представление; далее применяются одномерные свертки по временному и канальному измерениям, а также MLP-голова для генерации итоговых эмбеддингов. \textit{Справа:} энкодер фМРТ включает персонализированное маскирование мозга, векторизацию вокселей и проекцию в латентное пространство через индивидуальный линейный слой; завершающая обработка осуществляется с использованием MLP с остаточными соединениями}
    \label{fig:encoders}
\end{figure}

\paragraph{Эмбеддинги изображений.}
Для извлечения визуальных эмбеддингов используется предобученная модель CLIP ViT-H/14~\cite{radford2021learning}, которая генерирует эмбеддинги изображений $\mathbf{Z}_I \in \mathbb{R}^{B \times d}$. Данная модель обеспечивает высококачественные, семантически насыщенные и хорошо обобщающиеся эмбеддинги, а также совместима с множеством генеративных моделей, предобученных на CLIP-пространстве.

\paragraph{Контрастивное обучение энкодера сигналов мозга.}
Основной целью контрастивного обучения является сопоставление нейронных представлений фМРТ--ЭЭГ с визуальными эмбеддингами, извлеченными с помощью замороженного энкодера CLIP, как показано на рис.~\ref{fig:scheme}. Обучение происходит на триплетах (фМРТ, ЭЭГ, изображение) и включает оптимизацию трех ключевых компонентов: энкодера фМРТ $\mathcal{E}_f$, энкодера ЭЭГ $\mathcal{E}_e$ и модуля объединения $\mathcal{F}$. Для вычисления логитов близости используется косинусное расстояние между нормированными эмбеддингами:
\begin{equation}\label{eq:norm}
    \mathbf{Z}_I^n = \frac{\mathbf{Z}_I}{\|\mathbf{Z}_I\|}, \quad \mathbf{Z}_c^n = \frac{\mathbf{Z}_c}{\|\mathbf{Z}_c\|},
\end{equation}
где $\mathbf{Z}_I$~--- эмбеддинги изображений, полученные с помощью замороженного CLIP, а $\mathbf{Z}_c$~--- объединенные эмбеддинги фМРТ--ЭЭГ. Матрица логитов $\mathbf{L}$ вычисляется как
\begin{equation}\label{eq:logits}
    \mathbf{L} = \mathbf{Z}_I^n \left(\mathbf{Z}_c^n\right)^{\top} \cdot e^{\tau},
\end{equation}
где $\tau$~--- обучаемый параметр температуры. В процессе контрастивного обучения рассматриваются все возможные пары <<объединенный эмбеддинг мозга~--- эмбеддинг изображения>> для максимизации эффективности обучения. Критическим аспектом данного подхода является наличие большого числа негативных примеров, что позволяет модели не только корректно сопоставлять положительные пары, но и четко различать негативные. Как показано на рис.~\ref{fig:contrastive_learning}, при размере батча $B$ формируется матрица логитов размером $B \times B$, где лишь $B$ пар, расположенных на главной диагонали, являются положительными соответствиями.

Функция потерь основывается на симметричной контрастивной оптимизации, аналогичной подходу CLIP~\cite{radford2021learning}, и представляет собой усредненную сумму двух кросс-энтропийных потерь:
\begin{equation}\label{eq:contrastive_loss}
    \mathcal{L}_{\text{contr}} = -\frac{1}{2B} \sum_{b=0}^{B-1} \left[ \log \left( \frac{\exp(\mathbf{L}_{bb})}{\sum_{j=0}^{B-1} \exp(\mathbf{L}_{bj})} \right) + \log \left( \frac{\exp(\mathbf{L}_{bb})}{\sum_{i=0}^{B-1} \exp(\mathbf{L}_{ib})} \right) \right],
\end{equation}
где $\mathbf{L}_{ij}$ обозначает элемент матрицы логитов $\mathbf{L}$ на позиции $(i, j)$.

\begin{figure}[h]
    \centering
    \includegraphics[width=0.85\textwidth]{figure_3.pdf}
    \caption{Иллюстрация контрастивного обучения и формирования батчей. Отрицательные примеры формируются автоматически. Косинусные расстояния между эмбеддингами замороженного CLIP и нейронными эмбеддингами используются в качестве логитов}
    \label{fig:contrastive_learning}
\end{figure}

\subsection{Процесс реконструкции стимулов}\label{sec3:gen}

Вводится совместное распределение изображений $\mathbf{x}$ и их скрытых визуальных представлений $\mathbf{z}_I$ при условии комбинированных нейронных эмбеддингов $\mathbf{z}_c$:
\begin{equation*}
p(\mathbf{x}, \mathbf{z}_I \mid \mathbf{z}_c) = p(\mathbf{z}_I \mid \mathbf{z}_c)\, p(\mathbf{x} \mid \mathbf{z}_I).
\end{equation*}
Здесь $p(\mathbf{z}_I \mid \mathbf{z}_c)$ моделирует согласование комбинированного представления фМРТ--ЭЭГ с высокоуровневым визуальным пространством CLIP (стадия~I), а $p(\mathbf{x} \mid \mathbf{z}_I)$ описывает генерацию изображения при условии уточненного визуального эмбеддинга (стадия~II).

\paragraph{Стадия I: диффузионная модель в пространстве CLIP.}
В пространстве визуальных эмбеддингов CLIP обучается легкая диффузионная модель вида
\begin{equation*}
\boldsymbol{\epsilon}_{\boldsymbol{\theta}}(\mathbf{z}_I^{(t)},\, t,\, \mathbf{z}_c),
\end{equation*}
где $\mathbf{z}_I^{(t)}$~--- зашумленная версия истинного эмбеддинга $\mathbf{z}_I$ на шаге $t$ прямого диффузионного процесса, а $\mathbf{z}_c$~--- комбинированный эмбеддинг, полученный модулем слияния. Обучение осуществляется по методу DDPM (англ. Denoising Diffusion Probabilistic Models, вероятностные диффузионные модели расшумления)~\cite{ho2020denoisingdiffusionprobabilisticmodels} с функцией потерь
\begin{equation}\label{eq:prior_loss}
\mathcal{L}_{\text{prior}}(\boldsymbol{\theta}) = \mathbb{E}_{\mathbf{z}_I, \mathbf{z}_c, t, \boldsymbol{\epsilon}} \left\| \boldsymbol{\epsilon} - \boldsymbol{\epsilon}_{\boldsymbol{\theta}}(\mathbf{z}_I^{(t)}, t, \mathbf{z}_c)\right\|_2^2.
\end{equation}
На инференсе априорная диффузионная модель генерирует оценку $\widehat{\mathbf{z}}_I \approx \mathbf{z}_I$ итеративной процедурой расшумления, согласующей $\mathbf{z}_c$ с визуальным пространством CLIP. Кодировщики $\mathcal{E}_f$, $\mathcal{E}_e$ и модуль слияния $\mathcal{F}$ остаются замороженными, что упрощает схему и стабилизирует обучение.

\paragraph{Стадия II: генерация изображения при условии $\widehat{\mathbf{z}}_I$.}
Для реконструкции изображения используется предобученная латентно-диффузионная модель SDXL~\cite{podell2023sdxlimprovinglatentdiffusion}:
\begin{equation*}
\boldsymbol{\epsilon}_{\boldsymbol{\phi}}(\mathbf{x}^{(t)}, t, h(\widehat{\mathbf{z}}_I)),
\end{equation*}
где $\mathbf{x}^{(t)}$~--- зашумленный латент изображения, а $h(\cdot)$~--- проекционный адаптер IP-Adapter~\cite{ye2023ipadaptertextcompatibleimage}, переводящий $\widehat{\mathbf{z}}_I$ в формат условия, совместимый с SDXL. Модель $\boldsymbol{\phi}$ не дообучается.

\paragraph{Итоговый процесс реконструкции.}
\begin{enumerate}
    \item Контрастивное обучение выравнивает $\mathbf{z}_c$ с $\mathbf{z}_I$ (см. раздел~\ref{sec3:contrastive}).
    \item Априорная диффузионная модель генерирует согласованное визуальное представление $\widehat{\mathbf{z}}_I$ по $\mathbf{z}_c$.
    \item Предобученная SDXL, обусловленная через $h(\widehat{\mathbf{z}}_I)$, синтезирует изображение $\widehat{\mathbf{x}}$.
\end{enumerate}
